\documentclass{article}
\usepackage{hyperref}

\title{User Manual\\Final Project}
\author{Elvis Montiel, David Montiel, Steven Mayorquin, Vincent Lozano}
\date{Snapshot 3}

\begin{document}
\maketitle
\tableofcontents
\newpage

\section{Practice Simulation Disclaimer}

This repository is only a practice simulation for a final project. It is based on the senior design project Want2Remember. This repository is not affiliated with the original project, original design team, Want2Remember, We2Link, or any related organization. No actual application is being developed or distributed here.

\section{Project Overview}

This simulated project is a memory-support application for users who need help organizing memories, contacts, appointments, reminders, and daily tasks.

\section{Project Goals}

The goal of this project simulation is to demonstrate planning, documentation, project management, testing, Docker planning, and GitHub snapshot development.

\section{Jira Link}

Jira board: \url{https://ellymonty12.atlassian.net/jira/software/projects/CS38/boards/36}

\section{GitHub Repository}

Repository: \url{https://github.com/ellymonty12-web/cs3338-final-repo.git}

\section{How to Access}

This is not a working application. To access the project materials, open the GitHub repository and review the documentation folders.

\section{Snapshot 1 Status}

Snapshot 1 includes the repository setup, initial SDD, initial SRS, Docker Compose plan, Jira planning, and initial TestRail documents.

\section{Snapshot 2 User Instructions: Creating a Memory}

In the simulated application, a user would create a memory by logging in and opening the home dashboard. The user would select the Create Memory option, choose a memory template, fill in the requested information, and save the memory.

\subsection{Template Options}

The planned template options include appointment, medication, task, contact-related memory, password entry, and custom note. These templates are intended to make memory entry more organized and easier for users who may need memory support.

\subsection{Important Note}

This repository does not contain a working memory creation feature. This section describes the simulated user flow for the class project.

\section{Snapshot 3 User Instructions: Contacts and Search}

Snapshot 3 adds simulated user instructions for contact management and search/filter functionality.

\subsection{Viewing Contacts}

In the simulated application, a user would open the home dashboard and select the Contacts option. The Contacts page would display saved people connected to the user's memories.

\subsection{Viewing Contact Details}

The user would select a contact from the Contacts page to view more information. The Contact Details page may show the contact's name, relationship, phone number, email address, notes, and related memories.

\subsection{Searching Memories}

The user would open the Search page or use the search bar from the memory list. The user could enter a keyword to search saved memories.

\subsection{Filtering Memories}

The user could filter memories by category, date, contact, or template type. This would help the user locate information more quickly.

\subsection{Important Note}

This repository does not contain a working contacts or search feature. This section describes the simulated user flow for the class project.

\end{document}